\section{Fundamentos y estado del arte}
\label{sec:arte}
Esta sección recorre los conceptos necesarios para el resto del trabajo: qué se entiende por calidad de una imagen y, en particular, de una imagen facial; qué son los modelos de visión-lenguaje; y cómo se han empezado a usar para evaluar calidad.

\subsection{Evaluación de calidad de imagen y calidad biométrica facial}
\label{sec:iqa}
\pendiente{Definir brevemente la evaluación de calidad de imagen (\textit{Image Quality Assessment}, IQA) y los factores de degradación más comunes: iluminación, desenfoque, ruido, resolución y compresión.}

La forma más fiable de medir la calidad de una imagen es preguntar a un grupo de personas, pero es cara y lenta. Por eso la IQA objetiva intenta predecir esas opiniones con modelos automáticos, que se entrenan o se validan con la puntuación media que dan los observadores (\textit{Mean Opinion Score}, MOS)~\cite{Zhang_2025_TOMM}. Estos métodos se agrupan según la información de referencia que necesitan. Los de referencia completa (FR-IQA) comparan la imagen con una versión original sin degradar, como hace SSIM~\cite{Wang_2004_TIP}. Los de referencia reducida (RR-IQA) solo necesitan algunas características de esa original. Y los que no usan referencia (NR-IQA) trabajan únicamente con la imagen que se quiere evaluar, como BRISQUE~\cite{Mittal_2012_TIP}. En biometría casi nunca existe una imagen original con la que comparar, así que la FIQA es, en la práctica, un problema sin referencia.

\paragraph{Calidad perceptual frente a utilidad biométrica.} En biometría, una imagen facial tiene buena calidad si ayuda al sistema a reconocer a la persona sin equivocarse~\cite{Schlett_2022_CSUR}. Esto no siempre coincide con lo que a una persona le parece una buena foto. Una imagen puede ser nítida y estar bien iluminada y, aun así, ser poco útil si la cara aparece de perfil o tapada en parte. Para comprobar si un método de calidad mide utilidad, se suele observar cómo baja el error del reconocedor al descartar las imágenes que el método considera peores. Si el error baja deprisa, el método está acertando.

Hay trabajos que también se presentan como FIQA pero que miden calidad perceptual. Es el caso de DSL-FIQA~\cite{Chen_2024_CVPR}, un modelo basado en transformers que predice el MOS de retratos y que se fija especialmente en zonas de la cara como los ojos y la boca. Está pensado para tareas como la restauración de imágenes, y sus propios autores separan esta evaluación de la reconocibilidad de la cara. El VQualA 2025 FIQA Challenge~\cite{Ma_2025_ICCVW} sigue la misma línea perceptual.

\paragraph{Métodos FIQA especializados.} Los métodos pensados para la utilidad biométrica suelen apoyarse en el propio modelo de reconocimiento. SER-FIQ~\cite{Terhorst_2020_CVPR} genera varias representaciones (\textit{embeddings}) de la misma imagen activando y desactivando partes de la red al azar, y considera que la imagen es de buena calidad si esas representaciones se parecen entre sí. SDD-FIQA~\cite{Ou_2021_CVPR} compara cuánto se parece una imagen a otras de la misma persona y cuánto a las de personas distintas; cuanto más separadas estén ambas cosas, más útil es la imagen, y con esa medida entrena una red que predice la calidad de imágenes nuevas.
\pendiente{Añadir una o dos frases sobre FaceQAN~\cite{Babnik_2022_ICPR}, eDifFIQA~\cite{Babnik_2024_TBIOM} y ViT-FIQA~\cite{Atzori_2025_ICCVW}, que son los métodos con los que se comparan los VLMs en la sección~\ref{sec:analisis}.}

Estos métodos funcionan bien, pero se comportan como cajas negras: devuelven una puntuación sin decir qué ha fallado ni cómo corregirlo~\cite{Saritas_2026_FG}. En un control de fronteras o al hacerse la foto del pasaporte, con la persona delante de la cámara, eso es un problema práctico, porque el sistema no puede pedirle que se acerque o que mire de frente si no sabe cuál es la causa. Schlett et al.~\cite{Schlett_2022_CSUR} ya señalaban como reto que la FIQA estime la utilidad y, a la vez, explique por qué una imagen es mala. Esta es la principal motivación para probar con VLMs.

\subsection{Modelos de visión-lenguaje: fundamentos y arquitecturas}
\label{sec:vlms}
\pendiente{Explicar de forma divulgativa qué es un VLM y cómo relaciona imagen y texto. Presentar los dos tipos principales con un ejemplo de cada uno: los contrastivos, como CLIP~\cite{Radford_2021_ICML}, y los generativos, como LLaVA~\cite{Liu_2023_NeurIPS} y Qwen-VL~\cite{Bai_2025_arXiv}.}

Hay varias formas de usar un VLM para una tarea nueva. En \textit{zero-shot} el modelo se usa tal cual, sin entrenamiento adicional: basta con describirle la tarea en una instrucción de texto, o \textit{prompt}. En \textit{few-shot} se le muestran además unos pocos ejemplos resueltos. Con el ajuste fino (\textit{fine-tuning}) se vuelve a entrenar parte del modelo con datos de la tarea. El \textit{zero-shot} es el más barato, pero el resultado depende mucho de cómo se redacte el prompt.

Conviene distinguir también entre VLMs generalistas, entrenados para entender imágenes de cualquier tipo, y modelos adaptados a caras, como Face-LLaVA~\cite{Chaubey_2026_WACV} o FaceLLM~\cite{Shahreza_2025_ICCVW}. Por último, un VLM no tiene por qué trabajar solo. En un sistema híbrido, un algoritmo especializado se encarga de medir y el modelo de lenguaje se encarga de explicar el resultado. Los dos enfoques, el VLM generalista en \textit{zero-shot} y el sistema híbrido, se analizan en la sección~\ref{sec:analisis}.

\subsection{Aplicación de los VLM a la evaluación de calidad de imagen facial}
\label{sec:vlm-calidad}
El survey de Zhang et al.~\cite{Zhang_2025_TOMM} ofrece una clasificación útil de los trabajos que usan grandes modelos para evaluar calidad. Por un lado está la evaluación mediante prompts: se le pide al modelo que valore la imagen, a veces describiéndola antes de puntuarla, o que compare varias imágenes. Por otro está la evaluación basada en características, que incluye los métodos construidos sobre CLIP, los modelos multimodales usados en \textit{zero-shot} y los modelos ajustados específicamente para calidad, como Q-Align~\cite{Wu_2024_ICML} o DepictQA~\cite{You_2024_ECCV}.

En IQA general los resultados son prometedores pero todavía modestos. En la comparación que recoge el survey, el mejor modelo en \textit{zero-shot} alcanza una correlación con la opinión humana de alrededor de 0,6, frente a unos 0,35 de CLIP. Lo más relevante para este trabajo es que las caras apenas aparecen en el survey, y cuando lo hacen se tratan como retratos que se valoran por su aspecto. La utilidad biométrica no se menciona. A nuestro juicio, este vacío es el que justifica los trabajos específicos que se analizan en la siguiente sección.
\pendiente{Resumir en pocas frases Q-Align~\cite{Wu_2024_ICML} (IQA general) y CLIB-FIQA~\cite{Ou_2024_CVPR} (FIQA con CLIP), separando la evidencia de calidad general de la facial.}
